\documentclass[letter,12pt]{article}

\usepackage[T1]{fontenc}
\usepackage{lmodern}
\usepackage{textcomp}
\renewcommand*\familydefault{\sfdefault}

\usepackage[spanish]{babel}
\usepackage[utf8x]{inputenc}

\usepackage[pdftex]{graphicx}
\usepackage{pifont}
\usepackage[
pdfauthor={Carlos Caballero Burgoa},%
pdftitle={La unificación},%
colorlinks,%
citecolor=black,%
filecolor=black,%
linkcolor=black,%
%urlcolor=black
pdftex]{hyperref}

\usepackage{fancyhdr}
\usepackage{lastpage}
\pagestyle{fancy}

% Para la primera página
\fancypagestyle{plain}{
\fancyhead[l]{}
\fancyhead[r]{}
\fancyhead[c]{}
\renewcommand{\headrulewidth}{0.5pt}
\fancyfoot[l]{SCESI \\ Sociedad Científica de Estudiantes de Sistemas e Informática}
\fancyfoot[c]{}
\fancyfoot[r]{\thepage/\pageref{LastPage}}
\renewcommand{\footrulewidth}{0.5pt}}

% Para el resto de páginas
\lhead{La unificación}
\chead{}
\rhead{\includegraphics[width=0.1\textwidth]{scesi.png}}
\renewcommand{\headrulewidth}{0.4pt}
\lfoot{SCESI \\ Sociedad Científica de Estudiantes de Sistemas e Informática \\ 
\url {http://scesi.fcyt.umss.edu.bo}}
\cfoot{}
\rfoot{\thepage/\pageref{LastPage}}
\renewcommand{\footrulewidth}{0.4pt}

\title{\bf La unificación}
\author{Carlos Caballero Burgoa} 

\begin{document}
\maketitle
\begin{center}\includegraphics[width=0.48\textwidth]{yachay.png}\end{center}
\begin{center}\url {http://scesi.fcyt.umss.edu.bo}\end{center}
\pagebreak

\tableofcontents
\pagebreak

\section{Introducción}
Este documento trata los temas referentes a las posibilidades de crecimiento de los servicios ofrecidos por la sociedad científica, y su necesidad de unificación; para conseguir el valor sinérgico capaz de fortalecer los objetivos y visión de tal institución.

\section{Antecedentes}
Con la llegada de nueva infraestructura de hardware, y como consecuencia directa de esta; la sociedad científica tiene ahora la posibilidad de proveer a sus proyectos y grupos, el espacio y
las funcionalidades que estas necesiten.

A continuación se detallan los sitios que están (o deberían estar) en funcionamiento actualmente:

\begin{itemize}
\item babel.scesi.memi.umss.edu.bo
\item repoman.scesi.memi.umss.edu.bo
\item wiki.scesi.memi.umss.edu.bo
\item awesong.scesi.memi.umss.edu.bo
\item gallery.scesi.memi.umss.edu.bo
\item serendipia.scesi.memi.umss.edu.bo
\item nonchalant.scesi.memi.umss.edu.bo
\end{itemize}

Incluidos estos sitios, se percibe una tendencia creciente al uso de otros dominios, que estarán 
enfocadas a objetivos muy puntuales, y cuyas funciones pueden caer fácilmente en redundancia.

\section{Definición del Problema}
La sociedad científica de sistemas e informática ha venido proveyendo una creciente variedad de servicios; ya sea desarrollando con fines académicos muchos de ellos, o implantando
soluciones conocidas de software libre. Esta amplia gamma de aplicaciones ha creado un problema en lo que respecta a procedimientos, a menudo repetitivos sobre cada aplicación.

Con la llegada de nuevos integrantes y el fomento hacia el uso y desarrollo de software libre, que
es meta de esta sociedad; el problema antes mencionado comenzará a ser critico en lo referente a
interconectividad de estos proyectos.

Una de las barreras para solucionar esta potencial situación será el grado de conocimiento que se
debe tener respecto a los proyectos y sus formas de interrelación. Tal barrera podría ser superada
aprovechando el uso de la plataforma \emph{yachay}, debido en gran parte a la afinidad y facilidad
de comunicación con los creadores de tal herramienta.

Por lo mencionado anteriormente se define el problema como:

La gran variedad de software potencial utilizado y por utilizar en la sociedad científica conducirá 
en breve a un conjunto muy dispar de servicios inconexos.

\section{Objetivo General}
Implantar sobre la infraestructura de la sociedad científica, la plataforma \emph{yachay} para
unificar los servicios provistos por tal institución.

\section{Objetivos Específicos}
\begin{itemize}
\item Crear un marco de trabajo en el que puedan ser captadas todas las iniciativas de tipo local para ampliar el numero de estas.
\item Aprovechar el conocimiento acerca de la herramienta, para proveer funciones especializadas.
\item Unificar todos los servicios que ofrece actualmente y los que ofrecerá en un futuro la
sociedad científica para mejorar la difusión y por ende la utilidad de estas.
\item Proveer de funcionalidades comunes de modo que estas no sean una dificultad para los objetivos de los proyectos.
\end{itemize}

\section{Herramientas}
Para este proyecto, se requerirá utilizar la plataforma \emph{yachay}, como dominio principal de
la sociedad científica, partiendo de ahí, sindicar los proyectos de desarrollo de software al
sitio.

A continuación se detallan los pasos a seguir:

\begin{enumerate}
\item Puesta en marcha de la plataforma \emph{yachay} sobre el dominio principal de la sociedad
científica.
\item Creación de roles específicos al dominio de la sociedad científica.
\item Creación de cuentas para todos los miembros de la sociedad científica.
\item Utilización del dominio principal para autentificación sobre todos los proyectos de la
sociedad científica que requieran manejo de cuentas, basado en la tecnología OpenID.
\item Reutilización de los espacios virtuales de la plataforma \emph{yachay} para complementar funcionalidades comunes de los proyectos (comentarios, etiquetado, rankeo, blogging, rss, boletines, etc.).
\end{enumerate}

\section{Justificación}
El tener un sitio unificado, que englobe todos los servicios será en ultima instancia el factor
determinante sobre el que se apoyan muchas de las premisas que son objetivos de la sociedad
científica; tanto la difusión y apoyo al modelo de conocimiento libre, como la promoción de la
investigación científica, en este momento parecen carencias por omisión cuyo resolución debería ser
prioritaria.

Con cada idea de la sociedad científica convertida en un proyecto, recaen sobre sus integrantes un
conjunto de responsabilidades; debido a la falta de una metodología formal, la situación actual
requiere invertir mucho mas tiempo en actividades de organización que las que podrían conseguirse
con la estandarización de muchos de los procedimientos usados en proyectos mas desarrollados.

\end{document}          
